\chapter{问题、求解尝试与仍然缺少的环节}
\section{问题一：离子先热还是电子先热，为什么同一最终跳跃仍不够}
\subsection{已有尝试：共速双温慢激波}
Longcope与Bradshaw先取偏斜磁场重联后产生的流动，再以一维双温碰撞流体计算激波和传导前沿。此问题无电流的特殊约化使两物种共用密度和速度，电子惯性被忽略，电子压力梯度由电场平衡；它没有自行求解Hall扩散区。离子黏性把流能变成离子内能，库仑交换把能量送给电子，电子热传导又影响激波上游。\citep{L164}

共速、近各向同性时，可将两温方程示意写成
\begin{align}
\tfrac32nk_BD\theta_i+p_i\nabla\cdot\vct u
 &=\vct\tau_i:\nabla\vct u-\nabla\cdot\vct q_i+C(\theta_e-\theta_i),\\
\tfrac32nk_BD\theta_e+p_e\nabla\cdot\vct u
 &=-\nabla\cdot\vct q_e+C(\theta_i-\theta_e),\quad C>0.
\end{align}
这是展示各通道的简化式，不冒充原文所有系数。均匀无其他源时，温差$\delta\theta=\theta_i-\theta_e$满足$D\delta\theta=-4C\delta\theta/(3nk_B)$；因此温差松弛时间为$3nk_B/(4C)$。若文献以单一物种能量交换时间定义$C$，二者可能差因子2，比较时必须对齐定义。

原文计算分出离子黏性子激波、下游离子冷却区和电子传导前沿。早期离子温度可超出最终平衡温度，电子升温滞后；远下游平衡跳跃接近单温结果，却不代表中间路径相同。此结论超出H1的能力，但H2若使用同样两温、交换和输运方程，在适用极限内也应能够得到。因此这是对单温约束的修正，不是对全部Hall模型的否定。

\subsection{成立条件与缺口}
原文使用电子、离子Prandtl数约0.013、0.33，温度依赖的经典输运；稳态由射击法求解，另计算时变激波。部分强激波的早期离子热流接近或超过自由流量级，黏性修正甚至可能违反压力正定，恰好暴露经典离子闭合的边界。模型未包含X点直接电子加热，也未做实验两温的定量验证。不能用其电子热流较小的结果排除其他几何的直接电子加热。
\loc{L164}{本地45页版，式(4)--(7)、(14)--(18)；图4--7、9--14；物理页10、12、17、19--20、24、26、28、30、32、35、37。}

\section{问题二：双温Hall已有热点，完整双流体究竟增加什么}
\subsection{Browning与HIT-SI给出的基线}
Browning等的MAST磁环合并研究前列方程演化总压，后续图3才是初步独立两温扩展。离子黏性供热与电子电阻、超电阻供热产生不同空间图像，离子下游双峰与装置温升位置有联系。超电阻参数同时参与破冻结和电子热源，故热点并非仅由Hall项决定。本文未列该两温图专用网格和收敛序列，不可从总模型介绍补猜。\citep{L223}

HIT-SI的三维双温Hall模型更加完整地给出物种热源、交换和壁面热流。PSI-Tet与NIMROD中电子随电子流对流、压缩，欧姆热归入电子，黏性热归入离子；增强电子质量用于简化惯性项。模型能得到随注入频率变化的$T_i>T_e$，但装置功率包含注入器及反转过程，不能把全部热源归入单一X线。两代码比较还夹杂几何差别。\citep{Kaptanoglu2020}

\subsection{增加完整两流体后的可检验差别}
从H2到保留电子惯性的五矩，新增电子流能和快速电磁/电荷响应，但没有增加压力各向异性。若目标是远大于$\lambda_D,d_e$的下游温度，且两者采用一致热源和边界，这些新自由度可能只产生小修正；不能预设完整Maxwell一定改进加热。反之，在电子层，$D_e\vct u_e$和$\vct P_e$的非理想支撑可能改变电功位置。应通过局地欧姆项和能量项诊断确认哪一项造成差异。

Li与Livescu是很重要的反例：理论方程保留Braginskii物种输运，但实际四项验证强制$T_i=T_e$，删除独立离子能量演化。其六阶紧致差分、物理耗散和网格收敛说明数值框架的能力；并没有验证两物种加热分配。其所谓two-fluid又忽略电子惯性与位移电流，与通常“完整全Maxwell双流体”的定义不同。\citep{L247}
\loc{L223 / Kaptanoglu2020 / L247}{分别为图3（PDF7页）；式(1)--(13)、表I及图3--4（PDF3--4、8--9页）；第IV节关于$T_i=T_e$的实际测试约束与附录输运。}

\section{问题三：离子减速、压力支撑和温升是否由同一闭合解释}
\subsection{从局地松弛到非局地热流}
Wang等用五矩和十矩比较压力结构。十矩的有限压力松弛可以改善X点附近电子流及非对角压力，而强松弛近五矩会产生参照中没有的岛链。其图像常按相同重联磁通而非相同时间比较；$P_{xy}$峰值仍有显著误差。它证明新增张量自由度有用，并未证明温度、热流或全过程均正确。\citep{Wang2015}

Ng等2015年的Fadeev岛合并直接比较Hall和十矩。局地十矩改变离子压力支撑，却仍难得到参照中的大岛反弹与尺寸依赖。2017年非局地热流闭合使离子层厚从Hall的约$0.18d_i$改善到约$1.20d_i$，接近该算例参照的$1.15d_i$；但平均率的尺寸指数仍有偏差。\citep{Ng2015,Ng2017}

\subsection{温度梯度闭合的推进与未完成的加热账本}
Allmann-Rahn等2018年先检查热流散度，指出局地松弛的空间形状乃至符号可能错误，再以压力偏离各向同性的梯度闭合改善岛合并及反弹。其不同算例分别调节$k_0$，并非一套参数无校准成功。Ng等2020年改用温度梯度形式，平均重联率尺寸依赖和离子宽层进一步改善，并以同一典型$k_0d_s$测试重联与低混杂漂移不稳定性。\citep{L256,Ng2020}

这些研究对加热的意义在于：离子出流的减速长度、压力支撑和压力--应变交换位置会随闭合改变。\textbf{本报告的推论}是，仅用局地Hall温度峰去标定黏性，有可能掩盖错误的离子压力输运。原论文的主要成功指标却是层宽、反弹、率和不稳定性，而非经过实验验证的$Q_i^{\rm irr}/Q_e^{\rm irr}$。尚需把同一组计算中的完整压力功、热流、储能和边界项同时输出，才能闭合这一因果链。

成立条件包括二维岛合并几何、有限质量比和具体闭合波数。Ng2020指出电子各向异性结构仍不准确，跨场热流偏大；其附录还显示，对称稳态停滞点在零热流、无碰撞条件下可能趋向不物理压力。增加张量却继续设$Q=0$并不能保证稳态加热层可成立。
\loc{Wang2015 / Ng2017 / L256 / Ng2020}{Wang图4、8--10；Ng2017图1--5；L256图1--10与表1；Ng2020式(4)--(6)、图1--7、附录A。详细参数见附录。}

\section{问题四：电子各向异性有没有比标量温度更强的检验}
\subsection{可以进入流体方程的状态方程}
各向异性电子闭合将$p_{e\parallel},p_{e\perp}$写成局地$n,B$及外部来源参数的函数。这里的“局地函数”不意味着物理机制纯局地：外部热库和俘获/通过粒子的分界承担了非局地信息。新增的关键不是电子惯性，而是两个压力自由度和受俘获限制的热流响应。

Wetherton等2019年的MMS事件中，低密度区域电子接近等温Boltzmann响应，高密度区域更接近双绝热响应；一个固定标量电子压不能同时描述平行、垂直温度及转折。论文的数值对照使用H3D动理学离子加流体电子，必须归为混合模型；本报告仅将其用于验证可移植的电子闭合，不把离子结果算为双流体预测。\citep{Wetherton2019}

\subsection{从横向力平衡到有条件的温升标度}
该文强导场估计取峰值处$B\simeq B_g$，忽略电子压力对横向总压的主导作用，并取离子绝热$p_i\propto n^\Gamma$。由
\begin{equation}
p_{i\infty}\left[(n/n_\infty)^\Gamma-1\right]\simeq\frac{B_r^2}{2\mu_0}
\end{equation}
得$\widetilde n=(1+1/\beta_{ir})^{1/\Gamma}$，$\widetilde B=B_g/(B_r^2+B_g^2)^{1/2}$；$\beta_{ir}=2\mu_0n_\infty T_{i\infty}/B_r^2$。将二者代入电子状态方程，才得到标量温升。按原文低$\beta$渐近、$\Gamma=5/3$，
\begin{equation}
\frac{\Delta T_e}{m_iV_{Ar}^2}\simeq\frac{\pi}{36}\frac{T_{e\infty}}{T_{i\infty}}
\frac{B_r^2+B_g^2}{B_g^2}\beta_{ir}^{-1/5},
\quad T_{i\infty}\gg T_{e\infty},\quad B_g/B_r\gtrsim1.
\end{equation}
这里$m_iV_{Ar}^2=B_r^2/(\mu_0n_\infty)=2T_{i\infty}/\beta_{ir}$，因此指数由$-6/5$变为$-1/5$。系数$\pi/36$使用原文引用的渐近状态方程，本报告核验归一化及指数转换，不宣称重推其前序动理学积分。图2c的横轴是电子重联场beta，推导式中的$\beta_{ir}$是离子beta，不能混作同一参数。

原文典型范围得到约0.02的归一化电子温升；它是有条件的标度，不是电子占释放磁能的通用2\%，更不是不可逆熵产。电子温度的标度还借用了离子绝热闭合和峰值磁场假设，因此并非独立确定了离子加热。

2021年远下游研究进一步指出：局地磁俘获和有限通过时间造成的额外垂直增能，不能仅写成$n,B$的通用函数。它提出的能量修正需要轨道历史或额外状态信息；这正是有限局地闭合的可测边界。\citep{Wetherton2021}
\loc{Wetherton2019 / Wetherton2021}{前者PDF5--8页、图1--3与式(4)--(6)；后者PDF13--19页、表1、图6--10及式(4)。实验拟合条件见第4章。}

\section{问题五：出现三维不稳定性是否就解释了额外热化}
TenBarge等2019年以MMS 2015年10月16日事件参数计算全Maxwell五矩/十矩。十矩二维、三维的全域能量变化较近，离子内能增量较大；但电子各向异性峰虽接近观测幅度，峰值出现在错误的一侧。五矩还出现过强低混杂漂移活动与明显总能损失，不能把该模型的能量差直接解释为真实加热差。\citep{TenBarge2019MMS}

Allmann-Rahn等2021年及Bradshaw等2026年采用梯度型闭合，让十矩中的部分横向不稳定性得以增长。前者指出噪声幅度强烈影响启动与后续湍动；后者在类似MMS事件的三维非对称重联中得到局地闭合没有的漂移模、kink及通量绳。新增效果来自热流闭合对模态的阻尼选择，而非只来自“第三维”或“矩数更多”。\citep{E552,Closure2026}

Bradshaw等图6右还给出直接的反例约束：相对局地闭合，梯度闭合中的离子压力能和流能变化较小，电子压力能增长较快。恢复更多不稳定性并不自动带来更大的离子内能增量；这类全域能量曲线也没有分离不可逆熵产。

这两条线不能合并成一次独立重复：背景平衡、导场、质量比、闭合对象、数值扩散与种子均不完全相同。Bradshaw等的电子各向异性峰侧别仍不符，相关区域偏宽，且未给出完整三维网格/质量比收敛序列。\textbf{原文证明部分结构与模态出现，尚未证明其造成正确的两物种不可逆热化。}需要联合检验能量转移、热流、矩熵以及实验分布变化，而不能把“有湍流”直接当作耗散机制已知。
\loc{TenBarge2019MMS / E552 / Closure2026}{前者已读arXiv版图1及温度各向异性图；E552图4、6--9；Closure2026图1--11及附录A。正式版差异与收敛缺口保留在逐篇档案。}

\section{问题六：热流与多束有没有更合适的低成本表示}
\subsection{电子来源分群}
F635的两群电子十矩模型在固定离子背景上追踪电流片两侧来源。总热流由群内三阶矩、群间漂移和各群压力合成。它改善单群局地闭合的密度空洞与压力堆积，热流形状更接近参照，但幅度仍有偏差。这里“two-fluid”是两个电子群，\textbf{不是完整电子--离子双流体}；不能输出自洽离子加热。\citep{F635}

原文仍采用弱各向同性松弛；有限电子群数无法确保晚期每群保持单峰。图4在每个参照参数点取最佳时间相似度，不能读成两群流体在全参数扫描上的平均性能。图5色标不统一，也不能从图像近似断言幅度准确。其价值在于增加粒子来源历史这一低阶闭合缺失的信息；在本轮证据下，推广到离子参与、三维和实验仍属未验证。

\subsection{Landau/陀螺流体的时序检验}
L244从含电子两温、电子惯性和离子有限回旋半径的六场体系出发，但实际重联算例取冷离子子模型。Landau闭合降低平行温度扰动与增长率，使细尺度结构变平滑；按相同$\gamma t$比较的岛面积相似，不代表实际时间中的发展没有变慢。闭合还未表现出参照中的重联与加热延迟。\citep{L244}

这给出另一条超出H0/H1的途径：有限矩可近似部分平行相混，但自由能从解析变量消失，并不等于已在闭合中增加了相同数量的物种内能。若要称为加热预测，必须记录未解析自由能的去向或使用一致热库。实际冷离子算例没有回答离子、电子不可逆加热竞争；不能把理论六场变量清单当成已完成的双物种数值验证。

\subsection{为什么不按十三矩、二十矩的名称自动扩大结论}
独立演化三阶热流理论上能增加时间滞后和非局地响应，随后又需要四阶闭合、可实现性和稳定性控制。本轮检索发现相关会议摘要，但未取得足以核对离子--电子加热预算与实验的完整研究链，因此只列待核读，不将“矩数更多”写成已解决加热。Che的Buneman研究使用PIC诊断和近似流体分析，且算例未发展重联；它也不能替代独立流体重联验证。\citep{Che2014Buneman}
